\newpage
\section{Gleichstrommessbrücken}\label{sec:1}

\subsection{Abschätzung der Brückenwiderstände}\label{sec:1-1}
\subsubsection{Aufgabenstellung}

\begin{highlightblock}[Aufgabenbezeichnung]
\end{highlightblock}

\subsubsection{Verwendetes Material}

\begin{table}[h]
	\centering
	\caption{Verwendetes Material}
	\begin{tabularx}{\textwidth}{XXX}
		\textbf{Material} & \textbf{Verwendung} & \textbf{Anmerkungen}\\ \midrule
		& & \\
		& & \\
		& & \\
	\end{tabularx}	
\end{table}

\subsubsection{Versuchsaufbau}

\begin{figure}[h]
	\centering
	\begin{circuitikz}[scale=1.2]
		\draw (0,4) to[open, v={$U_H=\SI{3}{\volt}$}, o-o] (0,0)
			-- ++(2, 0) coordinate (k1)
			to[vR=$R_2$, i<_=$I_2$, *-] ++(0,2) coordinate (k2)
			to[R=$R_1$, i_<=$I_1$, -*] ++(0,2) coordinate (k3)
			-- ++(-2,0);
		\draw (k1) -- ++(3,0)
			to[vR=$R_4$, i<_=$I_4$] ++(0,2) coordinate (k4)
			to[vR=$R_3$, i_<=$I_3$] ++(0,2) -- ++(-3,0);
		\draw (k2) to[ni, v^=$U_0$, i>^=$I_0$, *-*] (k4);
		\node[yshift=-20pt] at ($(k2)!0.5!(k4)$) {$R_0$};
		\node[anchor=west] at ([shift={(0.5,-1)}]k4) {Abgleichwiderstand};
	\end{circuitikz}
	\caption{Wheatstone-Abgleichbrücke}
\end{figure}

\textbf{Geräte Einstellungen}

\begin{minipage}[t]{.48\textwidth}
	Voltmeter:
	\begin{itemize}[noitemsep]
		\item Messbereich: 
		\item Garantiefehlergrenzen: 
	\end{itemize}
\end{minipage}\hfill
\begin{minipage}[t]{.48\textwidth}
	Ohmmeter:
	\begin{itemize}[noitemsep]
		\item Messbereich: 
		\item Garantiefehlergrenzen: 
	\end{itemize}
\end{minipage}

\newpage
\subsubsection{Geräteparameter des Nullinstruments}

Geeignetes NI: Hoher Innenwiderstand, Garantiefehlergrenzen berücksichtigen
\begin{enumerate}[label=\alph*)]
	\item Kleinses Spannungsinkrement = Garantiefehler des Messfgerätes
	\item $R_4$ ist der Abgleichwiderstand, Logarithmisch Abgestuft.
		  Erhöhung von R4 um das Kleinste widerstandsinkrement soll das Kleinste Spannungsinkrement hervorrufen
	\item $R_1$ mit Ohmmeter messen. $\pm$ Garantiefehler
\end{enumerate}

\subsubsection{Abschätzung des Messobjekts}

Ergebnis für R1 plus Unsicherheit

\subsubsection{Berechnungen}

Gleichung 2.5 plus Abgleichbedingung

\begin{table}[h]
	\centering
	\caption{Abschätzung der Brückenwiderstände}\label{tab:brueckenabsch}
	\begin{tabularx}{\textwidth}{c|YYYY}
		Widerstandsverhältnis $R2:R1$ & $R_1 / \si{\ohm}$ & $R_2 / \si{\ohm}$ & $R_3 / \si{\ohm}$ & $R_4 / \si{\ohm}$ \\ \midrule
		$1:2$ \\
		$1:1$ \\
		$2:1$	
	\end{tabularx}
\end{table}

\subsubsection{Erkenntnisse}


